% ═══════════════════════════════════════════════════════════════════════════════
% METODOLOGÍA - INSTITUTO ATENEO DE COLIMA
% Estudio Documental: Corrupción en México (2006-2012)
% ═══════════════════════════════════════════════════════════════════════════════

\nota{
  Este capítulo detalla el diseño metodológico DIAC (CUAL + cuan) aplicado 
  conforme a Hernández Sampieri (2014). Se enfatiza el carácter documental 
  del estudio y los criterios de rigor aplicados.
}

\section{Enfoque y diseño de investigación (DIAC: CUAL + cuan)}

La investigación se desarrolló bajo un enfoque mixto formalizado como un Diseño Anidado Concurrente de Modelo Dominante (DIAC: CUAL + cuan), no experimental, retrospectivo y de estudio de casos múltiples \citep{hernandez2014}. En este diseño, ambas ramas de evidencia se recolectaron de forma simultánea (concurrente), pero el proyecto fue guiado por un método dominante de naturaleza cualitativa-documental (CUAL), dentro del cual se anidó un componente cuantitativo secundario de alcance descriptivo (cuan).

La integración metodológica operó de la siguiente manera: la rama dominante (CUAL) reconstruyó los procesos institucionales, las cadenas de evidencia jurídica y las categorías analíticas (captura estatal, opacidad contractual, dilución regulatoria y encuadre mediático) a partir de resoluciones judiciales, auditorías oficiales y doctrina especializada; de forma simultánea, la rama anidada (cuan) sistematizó indicadores estadísticos y financieros secundarios oficiales (puntajes del IPC 2012 y 2025, tasas de prevalencia, percepción y costo directo de la ENCIG 2025, montos desviados, porcentajes de sobrecosto, deuda subnacional y crecimiento de concesiones del IMSS). Ambas bases de datos se cruzaron en la fase de análisis e interpretación para generar metainferencias con consistencia interpretativa, permitiendo dimensionar numéricamente los daños institucionales explicados por la evidencia cualitativa \citep{hernandez2014}.

\section{Elección y justificación del tipo de investigación}

El estudio tuvo un carácter documental, descriptivo y analítico-explicativo, bajo un diseño no experimental, retrospectivo y de estudio de casos múltiples:

\begin{itemize}
  \item \textbf{Documental:} Se sustentó en la revisión sistemática y triangulada de fuentes primarias oficiales (auditorías de la ASF, expedientes del DOJ, comunicados de la UIF, sentencias de la SCJN, recomendaciones de la CNDH, informes de la CIDH y series del INEGI) y literatura académica mexicana e internacional dictaminada.
  
  \item \textbf{Descriptiva:} Detalló las propiedades, los actores, las cuantías financieras, las cronologías procesales y los indicadores nacionales que caracterizaron la corrupción en México entre 2006 y 2026.
  
  \item \textbf{Analítico-explicativa (en lógica cualitativa de casos):} Indagó cómo se configuró cada caso desde el punto de vista institucional, vinculando fallas en el diseño de contrapesos, planeación y regulación con sus consecuencias patrimoniales y en derechos humanos.
  
  \item \textbf{No experimental, retrospectiva y de estudio de casos múltiples:} No se manipularon variables ni se conformaron grupos experimentales; se examinaron hechos consumados entre 2006 y 2012 ---junto con sus derivaciones judiciales y fiscales hasta noviembre de 2026--- mediante el contraste estructurado de tres casos principales.
\end{itemize}

\section{Identificación de la población / universo y delimitación temporal}

El universo documental principal se delimitó al conjunto de expedientes de presunta corrupción, irregularidades administrativas graves y contrataciones públicas auditadas o judicializadas correspondientes a la Administración Pública Federal mexicana durante el sexenio 2006--2012, considerando sus resoluciones judiciales y pasivos financieros rastreables hasta noviembre de 2026. Como marco de referencia empírica nacional para la rama cuantitativa anidada, se emplearon los universos estadísticos oficiales de la ENCIG 2025 (población adulta residente en áreas urbanas de México medida por el INEGI) y del IPC 2025 (182 países evaluados por Transparencia Internacional).

Con el propósito de preservar el rigor de la delimitación temporal y evitar cualquier desviación del recorte 2006--2012, se estableció explícitamente que los dos episodios posteriores incluidos en el subtema 3.3 ---el saqueo hospitalario en Veracruz y la Refinería Olmeca (Dos Bocas)--- no formaron parte del universo nuclear del sexenio calderonista, sino que se incorporaron estrictamente como casos de contraste contemporáneo para evaluar la variable de encuadre mediático (framing).

\section{Tipo y procedimiento para muestras (casos principales y de contraste)}

Atendiendo a la naturaleza del estudio de casos múltiples, se aplicó un muestreo no probabilístico de tipo intencional (o dirigido por criterio teórico), estrategia adecuada para profundizar en fenómenos complejos cuando el objetivo es la riqueza documental y no la estimación estadística poblacional \citep{hernandez2014}. Del inventario documentado del sexenio 2006--2012 (que abarcó también Estela de Luz, ABB/CFE, BizJet, CPS penitenciarios, Odebrecht y Oceanografía), se seleccionó una muestra principal de tres casos emblemáticos (n = 3) ---Genaro García Luna, Refinería Bicentenario y Guardería ABC--- mediante tres criterios explícitos de inclusión:

\begin{enumerate}
  \item \textbf{Diversidad tipológica:} representar tres mecanismos analíticos diferenciados de corrupción (captura del aparato de seguridad pública, deficiencias de planeación y pasivos en megainfraestructura, y desregulación con conflicto de interés en servicios públicos subrogados).
  
  \item \textbf{Rigor probatorio oficial:} contar con evidencia emitida por órganos jurisdiccionales, fiscalizadores o de derechos humanos de ámbito nacional o internacional (cortes federales y civiles de EE. UU., UIF, ASF, SCJN, CNDH y CIDH).
  
  \item \textbf{Magnitud del impacto social y patrimonial:} documentar afectaciones verificables sobre la seguridad pública, el erario federal o estatal y derechos humanos fundamentales.
\end{enumerate}

Adicionalmente, para cumplir el cuarto objetivo específico relativo al encuadre mediático, se integraron por criterio teórico dos casos de contraste contemporáneo (Veracruz y Refinería Olmeca/Dos Bocas) cuya función metodológica se circunscribió a comparar la divergencia entre narrativas virales y evidencia oficial en materia de salud e infraestructura energética.

\section{Justificación y descripción de la técnica para recolección de datos}

Se empleó la técnica de análisis documental estructurado con triangulación de fuentes. Esta elección se justificó porque los esquemas de gran corrupción, lavado transnacional y captura institucional poseen un carácter clandestino que impide medirlos adecuadamente mediante cuestionarios de opinión local; en cambio, requieren reconstruir cadenas de evidencia cotejando resoluciones judiciales, auditorías técnicas, estadísticas oficiales y literatura académica. Para sistematizar la información se diseñaron y aplicaron dos instrumentos ---cuyos formatos modelo se presentan en los Apéndices A y B---:

\begin{itemize}
  \item \textbf{Ficha de Registro Jurídico-Institucional y Hemerográfico:} Instrumento cualitativo destinado a registrar el código de expediente, periodo, dependencias involucradas, tipología delictiva o administrativa, fuentes primarias trianguladas, hechos probados, posturas de descargo y categorías codificadas (Apéndice A).
  
  \item \textbf{Matriz Comparativa de Variables, Dimensiones e Indicadores por Caso:} Instrumento mixto de integración que cruzó las categorías cualitativas (mecanismo institucional, fallas de control y encuadre mediático) con los indicadores cuantitativos anidados (presupuestos, sobrecostos, deuda, sentencias pecuniarias y número de víctimas) en cada caso (Apéndice B).
\end{itemize}

\section{Método para análisis descriptivo y cualitativo-analítico}

En estricto apego a Hernández Sampieri et al. (2014, caps. 8 y 10), se declara expresamente que, debido a que la investigación empleó un muestreo no probabilístico intencional de tres casos documentales (n = 3), NO se aplicó estadística inferencial paramétrica ni no paramétrica para estimar parámetros poblacionales ni someter hipótesis a pruebas de significancia estadística. En su lugar, el procesamiento de la información integró dos procedimientos complementarios propios de un diseño DIAC (CUAL + cuan):

\begin{enumerate}
  \item \textbf{Estadística descriptiva (rama anidada cuan):} Se utilizó exclusivamente para organizar y describir los indicadores secundarios oficiales ---frecuencias absolutas y relativas, puntajes del IPC (2012 y 2025), tasas de prevalencia por cada 100 mil habitantes, promedios de gasto por víctima de la ENCIG 2025, porcentajes de sobrecosto presupuestal, evolución numérica de guarderías del IMSS y montos fijados en sentencias civiles--- sin proyectar inferencias probabilísticas fuera de sus fuentes originales.
  
  \item \textbf{Análisis cualitativo categorial, triangulación de fuentes e inferencia analítica comparada (rama dominante CUAL):} Constituyó el núcleo analítico de la investigación. Consistió en la codificación categorial de los documentos recopilados en torno a cuatro ejes (captura estatal, planeación presupuestal, regulación/conflicto de interés y framing mediático), la triangulación entre resoluciones judiciales, auditorías de la ASF y doctrina académica, y la inferencia analítica comparada entre los tres casos principales y los dos de contraste, con el objetivo de construir metainferencias trazables y analíticamente transferibles.
\end{enumerate}